\chapter{Phase 4: The Rust bridge}

\section{Goal}

Bring up the \texttt{hugo-bridge} bridge, which wires Asterisk's AMI events into the kernel's \texttt{uinput} subsystem, translating DTMF into virtual keystrokes.

This is the most software-heavy phase of the project. It is also used to document the bridge's internal architecture in detail.

\section{Code structure}

\begin{plaincode}
bridge/
|-- Cargo.toml                 # Crate manifest
|-- config.toml                # Runtime configuration
|-- hugo-bridge.service        # systemd unit
`-- src/
    |-- main.rs                # Entry point, tokio runtime
    |-- config.rs              # TOML parsing with serde
    |-- ami.rs                 # Minimal AMI client
    |-- keyboard.rs            # uinput virtual keyboard
    `-- watcher.rs             # File watcher (hot-reload)
\end{plaincode}

\section{Dependencies and design philosophy}

The bridge was designed to minimize external dependencies. The only \emph{crates} in \texttt{Cargo.toml} are:

\begin{itemize}
  \item \texttt{tokio}: async runtime (with specific \emph{features}, not \texttt{full}).
  \item \texttt{serde + toml}: configuration-file parsing.
  \item \texttt{tracing + tracing-subscriber}: structured logging.
  \item \texttt{clap}: command-line argument parsing.
  \item \texttt{uinput}: \emph{wrapper} over the \texttt{ioctl} API of \texttt{/dev/uinput}.
  \item \texttt{notify}: file watcher for hot-reload.
\end{itemize}

Conspicuously absent: any existing AMI client. The protocol is simple enough to implement by hand in about 150 lines, avoiding extra dependencies and keeping control over every byte that crosses the socket.

\section{The AMI protocol}

\emph{Asterisk Manager Interface} is a plain-text protocol over TCP. Its informal specification is:

\begin{itemize}
  \item TCP connection to \texttt{127.0.0.1:5038}.
  \item The server sends a banner: \texttt{Asterisk Call Manager/X.Y.Z\textbackslash r\textbackslash n}.
  \item Messages (in both directions) consist of \texttt{Key: Value} pairs separated by \texttt{\textbackslash r\textbackslash n}, terminated by a blank line (\texttt{\textbackslash r\textbackslash n\textbackslash r\textbackslash n}).
  \item The client must authenticate with an \texttt{Action: Login} message including \texttt{Username} and \texttt{Secret} fields.
  \item After that, the server sends asynchronous events whenever they occur.
\end{itemize}

A real-world message example:

\begin{plaincode}
Event: UserEvent
Privilege: user,all
UserEvent: HugoKey
Channel: PJSIP/102-00000001
Caller: 102
Key: 2
Uniqueid: 1715641234.1
(blank line)
\end{plaincode}

\section{Anatomy of the \texttt{ami.rs} module}

The AMI client is implemented in four functions:

\subsection{\texttt{run}: main loop with reconnection}

\begin{rustcode}
pub async fn run(mut cfg_rx: watch::Receiver<Arc<Config>>) -> Result<()> {
    let initial_cfg = cfg_rx.borrow().clone();
    let mut keyboard = VirtualKeyboard::new(
        &initial_cfg.keymap,
        initial_cfg.bridge.press_duration_ms,
    )?;

    let mut backoff_secs = 1u64;
    loop {
        let cfg = cfg_rx.borrow().clone();
        match connect_and_loop(&cfg, &mut keyboard, &mut cfg_rx).await {
            Ok(()) => {
                backoff_secs = 1;
            }
            Err(e) => {
                error!("AMI error: {}. Retrying in {}s...", e, backoff_secs);
                sleep(Duration::from_secs(backoff_secs)).await;
                backoff_secs = (backoff_secs * 2).min(30);
            }
        }
    }
}
\end{rustcode}

Three key decisions:

\begin{itemize}
  \item \textbf{The keyboard is created only once}, outside the connection loop. Recreating it between reconnections causes a brief window where the \texttt{/dev/input/eventXX} device disappears, which SDL2 handles poorly.
  \item \textbf{Exponential backoff} between retries (1, 2, 4, 8, 16, 30, 30, \ldots), capped at 30 seconds so we don't wait minutes when Asterisk restarts.
  \item \textbf{The config \emph{receiver} is passed by mutable reference} so the inner function can detect configuration changes while connected.
\end{itemize}

\subsection{\texttt{connect\_and\_loop}: connection, login and event loop}

After connecting the TCP socket, reading the banner and sending \texttt{Action: Login}, the main loop uses \texttt{tokio::select!} to multiplex two event sources:

\begin{rustcode}
loop {
    tokio::select! {
        msg = read_message(&mut reader) => {
            let msg = msg?;
            handle_event(&msg, &event_name, &key_field, keyboard).await;
        }
        changed = cfg_rx.changed() => {
            if changed.is_err() { return Ok(()); }
            let new_cfg = cfg_rx.borrow().clone();
            keyboard.update_keymap(&new_cfg.keymap);
            keyboard.set_press_duration(new_cfg.bridge.press_duration_ms);
            if new_cfg.ami != cfg.ami {
                return Ok(()); // forces reconnection
            }
        }
    }
}
\end{rustcode}

The \texttt{select!} blocks until one of two things happens. If a message arrives, it is processed. If the configuration changed (hot-reloaded by the watcher), it is applied to the keyboard, and if the AMI parameters changed, a reconnection is forced by returning \texttt{Ok(())} to the outer loop.

\subsection{\texttt{handle\_event}: filtering and action}

\begin{rustcode}
async fn handle_event(
    msg: &HashMap<String, String>,
    event_name: &str,
    key_field: &str,
    keyboard: &mut VirtualKeyboard,
) {
    if msg.get("Event").map(String::as_str) != Some("UserEvent") { return; }
    if msg.get("UserEvent").map(String::as_str) != Some(event_name) { return; }

    let Some(dtmf) = msg.get(key_field) else { return };
    let caller = msg.get("Caller").map(String::as_str).unwrap_or("?");

    info!("DTMF {} from {}", dtmf, caller);
    if let Err(e) = keyboard.press_dtmf(dtmf).await {
        error!("Failed to press key: {}", e);
    }
}
\end{rustcode}

Three successive filters discard AMI events the bridge doesn't care about. Only \texttt{UserEvent}s of type \texttt{HugoKey} (configurable) reach the key-injection path.

\subsection{\texttt{read\_message}: wire-format parser}

\begin{rustcode}
async fn read_message<R: tokio::io::AsyncBufRead + Unpin>(
    reader: &mut R,
) -> std::result::Result<HashMap<String, String>, BoxError> {
    let mut map = HashMap::new();
    let mut line = String::new();
    loop {
        line.clear();
        let n = reader.read_line(&mut line).await?;
        if n == 0 {
            return Err("AMI connection closed by server".into());
        }
        let trimmed = line.trim_end_matches(['\r', '\n']);
        if trimmed.is_empty() {
            if map.is_empty() { continue; }
            return Ok(map);
        }
        if let Some((k, v)) = trimmed.split_once(':') {
            map.insert(k.trim().to_string(), v.trim().to_string());
        }
    }
}
\end{rustcode}

Reads line by line until it finds a blank line, at which point it considers the message complete. Malformed lines (without a \texttt{:}) are silently discarded, typical behavior of tolerant parsers in legacy protocols.

\section{The \texttt{keyboard.rs} module}

The Linux kernel's \texttt{uinput} subsystem lets a \emph{user space} process create virtual input devices. The kernel exposes them just like any physical keyboard under \texttt{/dev/input/eventXX} and applications like DOSBox cannot tell them apart from real hardware.

\subsection{Creating the virtual keyboard}

\begin{rustcode}
let device = uinput::default()?
    .name("hugo-phone-virtual-keyboard")?
    .event(uinput::event::Keyboard::All)?
    .create()?;
\end{rustcode}

The constant \texttt{Keyboard::All} registers every key the kernel knows about. This matters for hot-reload: if we only registered the keys present in the initial \emph{keymap}, switching the mapping to a new key would require destroying and recreating the device, which interrupts the flow.

\subsection{Injecting a press}

\begin{rustcode}
pub async fn press_dtmf(&mut self, dtmf: &str) -> Result<()> {
    let Some(key) = self.keymap.get(dtmf) else {
        warn!("DTMF {:?} unmapped", dtmf);
        return Ok(());
    };
    self.device.press(key)?;
    self.device.synchronize()?;
    sleep(self.press_duration).await;
    self.device.release(key)?;
    self.device.synchronize()?;
    Ok(())
}
\end{rustcode}

Each press follows the cycle:

\begin{enumerate}
  \item Look up the symbolic key matching the received DTMF.
  \item Send a \texttt{press} event.
  \item Call \texttt{synchronize()} so the kernel emits a \texttt{SYN\_REPORT} and propagates the event to consumers.
  \item Wait \texttt{press\_duration} (typically 80\,ms).
  \item Send a \texttt{release} event and sync again.
\end{enumerate}

The pause between \emph{press} and \emph{release} is important: too short and games like Hugo don't detect the press; too long and the game's engine triggers the action in continuous repetition.

\section{The \texttt{watcher.rs} module}

The configuration \emph{hot-reload} lets you tweak the \emph{keymap} and the press duration without restarting the service. It is particularly useful during initial calibration.

\begin{rustcode}
let mut watcher = notify::recommended_watcher(move |res: notify::Result<Event>| {
    match res {
        Ok(evt) => { let _ = tx.send(evt); }
        Err(e) => error!("Watcher error: {}", e),
    }
})?;
watcher.watch(parent, RecursiveMode::NonRecursive)?;
\end{rustcode}

Two important details:

\begin{itemize}
  \item We watch the \textbf{parent directory}, not the file itself. Many editors save a file by writing to a temporary and renaming it on top of the original (\emph{atomic save}), which doesn't fire a \emph{modify} event on the original file but a \emph{create} in the directory.
  \item A 500\,ms \emph{debounce} is applied to avoid multiple reloads if the editor writes in several steps.
\end{itemize}

\section{Practical verification}

Test the bridge manually:

\begin{shellcode}
cd /opt/hugo-phone/bridge
sudo ./hugo-bridge --config config.toml
\end{shellcode}

Expected output:

\begin{plaincode}
INFO Config loaded from "/opt/hugo-phone/bridge/config.toml"
INFO Virtual keyboard created
INFO Watching ".../config.toml" for hot-reload
INFO Connecting to AMI on 127.0.0.1:5038 ...
INFO AMI banner: Asterisk Call Manager/9.0.0
INFO Authenticated to AMI as hugo-bridge
\end{plaincode}

Check that the virtual keyboard exists:

\begin{shellcode}
sudo apt install -y evtest
sudo evtest
\end{shellcode}

\texttt{hugo-phone-virtual-keyboard} must appear in the list. Select its number, then press DTMF keys in Linphone: evtest must show the corresponding press events.

\section{The complete flow with the softphone}

With the bridge running, DOSBox showing Hugo on the TV and Linphone registered:

\begin{enumerate}
  \item Dial \texttt{9000} from Linphone.
  \item Hear the beep in the earpiece.
  \item Press \texttt{2}. The character should jump/move up.
  \item Try \texttt{4}, \texttt{6}, \texttt{8}, \texttt{5}.
\end{enumerate}

In the bridge logs:

\begin{plaincode}
INFO DTMF 2 from 102
INFO DTMF 4 from 102
\end{plaincode}

\section{Calibration}

The most sensitive parameter is \texttt{press\_duration\_ms} in \texttt{config.toml}:

\begin{tomlcode}
[bridge]
press_duration_ms = 80
\end{tomlcode}

Symptoms and adjustments:

\begin{table}[h]
\centering
\begin{tabular}{ll}
\toprule
\textbf{Symptom} & \textbf{Adjustment} \\
\midrule
One press = several actions & Lower to 40--50\,ms \\
The press doesn't register & Raise to 120--150\,ms \\
Movement is choppy & Raise to 100\,ms \\
\bottomrule
\end{tabular}
\end{table}

Any change is applied on save (hot-reload), without restarting the service.

\section{Enable as a service}

Once validated manually:

\begin{shellcode}
sudo systemctl start hugo-bridge
sudo systemctl status hugo-bridge
sudo systemctl enable hugo-bridge
\end{shellcode}

Live logs:

\begin{shellcode}
sudo journalctl -u hugo-bridge -f
\end{shellcode}

\section{Troubleshooting}

\subsection{\texttt{Permission denied} on \texttt{/dev/uinput}}

\begin{itemize}
  \item The service runs as \texttt{root}, this shouldn't happen.
  \item If you're running it manually, use \texttt{sudo}.
  \item Check the udev rule at \texttt{/etc/udev/rules.d/99-hugo-uinput.rules}.
\end{itemize}

\subsection{The bridge connects but doesn't receive events}

\begin{itemize}
  \item Confirm that Linphone is calling \texttt{9000} and not some other extension.
  \item On the Asterisk console, does \texttt{UserEvent: HugoKey} appear when you press DTMF? If not, the problem is in the dialplan (Phase 2).
\end{itemize}

\subsection{DTMFs reach the bridge but Hugo doesn't respond}

\begin{itemize}
  \item DOSBox requires window focus. If it was launched over SSH with X-forwarding, it loses focus.
  \item Test the injection against another application: with the bridge running, open \texttt{nano} on the DOSBox TTY and verify that the cursor moves when you press DTMF.
\end{itemize}

\subsection{Cargo can't find dependencies during compilation}

\begin{itemize}
  \item Check Internet connectivity: \texttt{ping crates.io}.
  \item If DNS is restricted, configure:
  \begin{shellcode}
echo "nameserver 1.1.1.1" | sudo tee /etc/resolv.conf
  \end{shellcode}
\end{itemize}

\section{Phase wrap-up}

With this phase, the whole software flow is validated. All that remains is replacing the Linphone softphone with a real analog telephone connected to the HT801.
